% BHU PYQ -- Transcribed & Corrected
% Paper: CHESE31: Synthetic Methodologies (End-Term)
% Exam: Under Graduate Semester III Examination 2025-26 (NEP SEC)
% Department: Chemistry

\documentclass[11pt,a4paper]{article}
\usepackage[a4paper,top=2.5cm,bottom=2.5cm,left=2.8cm,right=2.8cm,headheight=14pt]{geometry}
\usepackage{amsmath,amssymb,chemformula}
\usepackage[T1]{fontenc}
\usepackage[utf8]{inputenc}
\usepackage{lmodern,microtype,fancyhdr,enumitem,hyperref,lastpage,parskip}

\hypersetup{
    colorlinks=true,
    urlcolor=black,
    linkcolor=black,
    pdftitle={CHESE31: Synthetic Methodologies -- BHU Sem III 2025-26 (NEP SEC)}
}

\pagestyle{fancy}
\fancyhf{}
\renewcommand{\headrulewidth}{0.4pt}
\renewcommand{\footrulewidth}{0.4pt}
\fancyhead[L]{\small\textit{Banaras Hindu University}}
\fancyhead[R]{\small\textit{CHESE31: Synthetic Methodologies (SEC)}}
\fancyfoot[C]{\small Page \thepage\ of \pageref{LastPage}}

\newlist{parts}{enumerate}{2}
\setlist[parts,1]{label=(\alph*),leftmargin=*,itemsep=4pt,topsep=2pt}
\newcommand{\pts}[1]{\hfill\textit{[#1]}}

\begin{document}

\begin{center}
  {\large\bfseries BANARAS HINDU UNIVERSITY}\\[2pt]
  {\normalsize Under Graduate Semester III Examination 2025-26 (NEP SEC)}\\[6pt]
  \rule{0.8\linewidth}{0.5pt}\\[6pt]
  {\large\bfseries SEC --- CHEMISTRY}\\[4pt]
  {\large\bfseries Paper No. CHESE31 : Synthetic Methodologies}\\[4pt]
  {\normalsize Time: 2:30 Hours\hfill Maximum Marks: 70}
\end{center}

\rule{\linewidth}{0.4pt}

\smallskip
\noindent(Write your Roll No. at the top immediately on the receipt of this question paper)

\smallskip
\noindent\textit{Note: Answer five questions including Question 1 which is compulsory. The figures in the right-hand margin indicate marks.}

\bigskip

\begin{enumerate}[label=\textbf{\arabic*.}]

    \item Answer all of the following:
    \begin{parts}
        \item What are the essential characteristics that define a `click' reaction?\pts{2}
        \item Write the products for the Ru-catalyzed azide--alkyne cycloaddition reaction.\pts{2}
        \item Write the product for the reductive amination reaction of ketone using $\ch{NaCNBH3}$.\pts{2}
        \item Write the names of missing reagents for selective functional group transformations.\pts{2}
        \item Why is the Grignard reagent prepared in an ether solvent?\pts{2}
        \item What is the Umpolung reaction? Explain it with suitable examples.\pts{2}
        \item Despite being highly hydrophobic, why do benzene or toluene still contain a substantial amount of dissolved water?\pts{2}
    \end{parts}

    \bigskip

    \item Answer the following:
    \begin{parts}
        \item Describe the Copper-Catalyzed Azide-Alkyne Cycloaddition (CuAAC) reaction. Why is CuAAC considered the most widely used click reaction?\pts{7}
        \item Write the thiol--ene click mechanism and explain its applications.\pts{7}
    \end{parts}

    \bigskip

    \item Answer the following:
    \begin{parts}
        \item What factors govern regioselectivity in Ruthenium-Catalyzed Azide-Alkyne Cycloaddition reactions?\pts{7}
        \item Define bio-orthogonal click chemistry and explain why it is important in chemical biology.\pts{7}
    \end{parts}

    \bigskip

    \item Answer the following:
    \begin{parts}
        \item Write the differences between $\ch{NaBH4}$ and $\ch{NaCNBH3}$ in terms of their reducing abilities and properties. Why is $\ch{NaCNBH3}$ a less powerful reducing agent than $\ch{NaBH4}$?\pts{6}
        \item Explain the organocatalyzed SuFEx click reaction and describe the mechanism in detail.\pts{4}
        \item Write the mechanism for the reduction of an ester and a nitrile using $\ch{LiAlH4}$.\pts{4}
    \end{parts}

    \bigskip

    \item Answer the following:
    \begin{parts}
        \item Write the product for the amino acid condensation reaction and explain the mechanism.\pts{5}
        \item How would you describe the Schlenk equilibrium?\pts{4}
        \item Describe how diphenylacetylene can be converted into the corresponding E- and Z-alkenes. Explain the conditions and the underlying reasoning for each transformation.\pts{5}
    \end{parts}

    \bigskip

    \item Answer the following:
    \begin{parts}
        \item Write short notes on liquid chromatography and adsorption chromatography.\pts{5}
        \item Write the mechanism for the reaction of carbonyl compounds with Grignard reagents, including the key transition states.\pts{5}
        \item Write the products with appropriate mechanism for: (i) Reaction of di-isopropyl ketone with isopropylmagnesium bromide, (ii) Reaction of excess benzaldehyde with phenylmagnesium bromide.\pts{4}
    \end{parts}

    \bigskip

    \item Answer the following:
    \begin{parts}
        \item Describe the origin of the blue color that appears when sodium or potassium metal dissolves in liquid ammonia. How does the electrical conductivity of the solution change with increasing amounts of dissolved metal?\pts{5}
        \item Write the mechanism of the Birch reduction and explain each step. What products are formed when naphthalene and anthracene undergo Birch reduction?\pts{5}
        \item Describe appropriate drying strategies for the following solvents: acetone, acetonitrile, diethyl ether, and N,N-dimethylformamide (DMF).\pts{4}
    \end{parts}

\end{enumerate}

\end{document}
